\chapter[Introduction]{Introduction}

\section{Contextualization}

For millennia, the essential value attached to mathematical proof has been its capacity to yield absolute certainty, establishing conclusions that can be universally accepted as true [p1, 1]. Securing the truth of a piece of knowledge and convincing others of its incontrovertibility is the ultimate ideal embodied by mathematical proof [p1, 2]. Under the traditional view, a mathematical claim is only considered fully justified when supported by a rigorous deductive proof, which allows human beings to transcend empirical observation and touch absolute certainty [p1, 3, 4].

The foundation of logical mathematical proof was laid in ancient Greece, moving away from purely visual or empirical demonstrations to rigorous deductive reasoning [p1, 9, 10]. Euclid’s Elements (c. 300 B.C.E.) served as the paramount exemplar of this method, establishing a system where propositions were derived entirely from a finite set of definitions, postulates, and axioms [p1, 11, 12]. For over two millennia, this axiomatic-deductive model functioned as the gold standard for verifying mathematical truth [p1, 12].

A profound paradigm shift occurred at the turn of the 19th century when mathematics evolved beyond explicit computation into highly abstract algebraic and set-theoretic methods championed by figures like Georg Cantor and Richard Dedekind [p1, 13]. The introduction of infinitary methods and the subsequent discovery of paradoxes in naive set theory sparked a "crisis of foundations" regarding whether these modern methods were actually consistent [p1, 13]. In response, David Hilbert launched his Beweistheorie (Proof Theory) in 1922 [p1, 14]. Hilbert’s program aimed to secure the foundations of mathematics by representing abstract reasoning through formal axiomatic systems governed by strict, mechanical rules of inference [p1, 14]. In this framework, mathematical proofs were stripped of intuition and treated as finite, combinatorial objects, allowing mathematicians to study the structure of mathematical reasoning itself [p1, 14, 15]. Although Kurt Gödel’s 1931 incompleteness theorems proved that no sufficiently complex system could prove its own consistency, Hilbert’s vision successfully demonstrated that all mathematical methods could, in principle, be formalized and reduced to a strict set of symbolic rules [p1, 16, 17].

Despite the theoretical possibility of perfect rigor, modern mathematics faces a severe human bottleneck. Manual proofs written in natural language are inherently prone to errors, and the mathematical literature is consequently filled with minor mistakes and misstatements that cause ongoing confusion [p1, 18, 19]. Even the most elite mathematicians are not immune to this; Andrew Wiles's celebrated 1993 proof of Fermat’s Last Theorem contained a critical gap that took a year of rigorous revision effort to fix [p1, 20, 21], and Fields Medalist Vladimir Voevodsky discovered a fatal error in one of his own major papers more than a decade after it was published [p1, 22]. As the boundaries of human knowledge expand, proofs are becoming incredibly long and overwhelmingly complex. This creates a massive burden on the peer-review process, as human reviewers share the same cognitive limitations and easily miss subtle logical gaps [p1, 23]. The current culture features "too-long proofs" that only a handful of competent referees in the world can understand, and these referees are often too busy writing their own lengthy proofs to properly scrutinize the work of others [p1, 24, 25]. For example, Fields Medalist Peter Scholze was deeply concerned that his foundational theorem on condensed mathematics might never be carefully checked by human referees due to its immense complexity [p1, 26].

Because human verification is so time-consuming and fragile, some provocateurs, like Doron Zeilberger, have even predicted a future where mathematicians might renounce absolute certainty altogether, accepting "almost certainty" generated by computers simply because rigorous proof has become too expensive and arduous to achieve [p1, 27, 28]. However, the mainstream mathematical community is instead turning to computer proof assistants to overcome human limitations. By painstakingly translating mathematics into machine-verifiable formal languages, mathematicians can bypass the bottleneck of human peer review and achieve absolute formal verification [p1, 23, 29].

The advent of digital computers in the second half of the 20th century profoundly transformed the mathematical landscape, making it natural to explore how machines could assist in representing, checking, or even producing mathematical proofs [p2, 1]. Early computer-assisted proofs were groundbreaking but fraught with practical and epistemic issues [p2, 2]. A famous early landmark was the 1976 proof of the Four Color Theorem by Kenneth Appel and Wolfgang Haken, which relied heavily on custom computer software to check the "reducibility" of 1,834 graphs [p2, 2]. However, this early method still required humans to painstakingly verify hundreds of pages of microfiche by hand to establish the "unavoidability" of these graphs, a tedious process that inadvertently introduced several minor, fixable errors [p2, 2]. Because these early computational proofs were partially opaque and relied on custom code, they were often accepted by the mathematical community only with reservation [p2, 3, 4]. Eventually, to resolve these doubts, mathematicians turned to formal proof assistants to engineer fully machine-checked proofs of these landmark theorems, including a complete formalization of the Four Color Theorem in Coq in 2005 and the Kepler conjecture in 2014 [p2, 4-8].

Automated Theorem Proving (ATP) systems are designed to take a system of axioms and find proofs autonomously [p2, 9, 10] by relying on traditional AI techniques like first-order theorem provers, SAT solvers, and SMT solvers. However, they face profound challenges with complex, creative leaps because they rely heavily on mechanical search strategies, causing a combinatorial explosion of the search space [p2, 11]. On the other hand, Interactive Theorem Proving (ITP) represents the modern standard for formal mathematics, with systems like Isabelle, Coq, and Lean leading the field [p2, 7, 15]. Unlike autonomous ATPs, ITPs operate on a paradigm of human-machine collaboration [p2, 16, 17]. In this setup, the human mathematician acts as the strategist, issuing commands—traditionally called "tactics"—that represent high-level logical inferences [p2, 16]. The ITP then applies these tactics, rigorously checks their syntactic and logical validity, and updates the explicit proof state, feeding a new list of subgoals back to the user [p2, 15, 16, 18]. This interactive, step-by-step verification ensures comprehensive logical validation, yielding what is sometimes called "super-human certainty" that a complex result is entirely free of hidden logical gaps [p2, 16, 19]. However, this absolute certainty comes at a high price. The major pain point of ITPs is that they require specialized coding knowledge, creating an incredibly steep learning curve for traditional mathematicians [p2, 16, 20]. Operating an ITP is structurally akin to software engineering: users must translate abstract mathematical concepts into a rigid, cryptic artificial language with strict syntactic rules and precise operational semantics [p2, 18, 21, 22].

Because the proof assistant enforces absolute precision, mathematicians cannot use the informal, narrative shortcuts common in pen-and-paper proofs [p2, 23, 24]. Instead, they must provide massive amounts of explicit logical detail, perfectly manage variable scopes, and precisely define the specific "type" of every mathematical object used (e.g., distinguishing between a natural number and a real number for an exponentiation operation) [p2, 23, 25]. This "translation gap" places a heavy cognitive burden on the user, making it difficult to formalize cutting-edge mathematics without extensive programming expertise [p2, 16, 26, 27].

The recent artificial intelligence boom has been overwhelmingly driven by the rapid scaling and advancement of Large Language Models (LLMs) [p3, 1]. These systems have transcended their initial purpose of basic text prediction, evolving into powerful engines capable of executing highly structured reasoning and complex symbolic tasks [p3, 1]. Historically, Natural Language Processing (NLP) focused heavily on conversational text and basic linguistic modeling. However, modern LLMs are developed through a foundational pretraining phase where they ingest massive, heterogeneous corpora of unlabeled text [p3, 2]. Because this training data includes not just everyday conversation but also textbook-style mathematical exposition, technical academic papers, online forum discussions, and billions of lines of source code, the models are forced to learn much more than basic grammar [p3, 2-4]. Through the simple objective of next-token prediction over this diverse data, NLP models have evolved to deeply internalize complex syntax, underlying semantics, world knowledge, and sophisticated discourse patterns [p3, 2].

To overcome the steep learning curve and specialized coding requirements of Interactive Theorem Provers (ITPs), researchers are increasingly applying Large Language Models (LLMs) to act as an automated bridge. By combining the vast pattern-recognition capabilities of neural networks with the absolute rigor of symbolic provers, AI aims to synthesize the best of both worlds: the intuitive, high-level strategic reasoning of humans and the flawless, low-level verification of machines [p4, 1], [p4, 2]. This integration is primarily driven by Autoformalization: the automatic translation of natural language mathematical statements and proofs into machine-verifiable formal code, such as Lean, Coq, or Isabelle [p4, 3], [p4, 4], [p4, 5]. The ultimate goal of autoformalization is to allow mathematicians to write in standard mathematical prose while an AI seamlessly translates their ideas into rigorous code behind the scenes [p4, 4], [p4, 6]. If fully realized, autoformalization could render vast amounts of existing mathematical literature programmable, vastly expanding the libraries available to automated provers and accelerating the discovery of new mathematics [p4, 5], [p4, 7].

However, autoformalization is an incredibly difficult task due to the profound clash between informal ambiguity and formal strictness. Traditional mathematical prose is embedded in a flexible narrative structure that relies heavily on approximations, linguistic conventions, and rhetorical shortcuts like "clearly" or "it follows that" [p4, 8], [p4, 9]. Human readers automatically reconstruct missing logical steps, resolve ambiguous references, and infer background context [p4, 3], [p4, 9]. In stark contrast, formal mathematics imposes an unforgiving symbolic discipline where every object must be perfectly typed, all implicit assumptions must be explicitly stated, and every inference must strictly follow predefined logical rules [p4, 10], [p4, 11]. A single incorrect tactic or minor deviation from the syntax can render an entire formal proof completely invalid [p4, 12]. Furthermore, natural language math suffers from "conceptual polymorphism," meaning the exact same informal term can map to entirely different formal definitions depending on the specific domain or level of abstraction [p4, 13], [p4, 14].

Because of this massive translation gap, naive LLMs exhibit significant structural shortcomings at complex formal mathematics without a symbolic grounding. LLMs operate probabilistically and lack a persistent, internal deductive "state" to track complex logical dependencies [p4, 15], [p4, 16], [p4, 17]. While their pretraining allows them to brilliantly mimic the textual footprints and discourse patterns of mathematical reasoning, they do not inherently verify logical truth [p4, 18], [p4, 19]. Consequently, when left to formalize mathematics autonomously, naive LLMs are highly susceptible to hallucinations—they frequently fabricate undefined mathematical concepts, invent non-existent library imports, misuse symbols, or mix syntax from different proof languages [p4, 20], [p4, 13]. They might generate arguments that look highly convincing and mathematically articulate on the surface, but are structurally riddled with logical fallacies, circular reasoning, and fatal gaps [p4, 21], [p4, 22]. To succeed, AI models cannot rely on language prediction alone; they must be symbolically grounded. This means they require tight integration with the proof assistant itself, relying on the formal compiler to act as an external working memory and strict verifier [p4, 23], [p4, 24]. By receiving step-by-step compiler error messages, syntax checks, and access to retrieved mathematical libraries, the AI can be grounded in actual logic rather than mere linguistic association [p4, 25], [p4, 26], [p4, 27].

To overcome the inherent limitations of pure generative AI, the near future of mathematical artificial intelligence is shifting decisively toward agentic systems and neuro-symbolic AI [p5, 1, 2]. Traditional large language models (LLMs) operate as probabilistic engines trained on next-token prediction; while they excel at mimicking the discourse patterns of human reasoning, they do not inherently maintain a persistent, verifiable deductive state [p5, 3-5]. Consequently, when tasked with complex mathematics, pure generative models are prone to hallucinating facts, missing subtle logical gaps, or producing arguments that appear plausible but are structurally invalid [p5, 6-8]. To solve this, researchers are embracing neuro-symbolic AI, a paradigm that synthesizes the complementary strengths of two distinct approaches [p5, 2]: On one side, machine learning (neural networks) provides vast pattern recognition capabilities, heuristic intuition, and the ability to process massive amounts of unstructured data [p5, 2]. On the other side, traditional symbolic AI provides rigor, relying on precise, mathematically specified semantics and strict rules of inference [p5, 2, 9, 10]. By bridging these domains, neural networks can propose intuitive, high-level strategic leaps, while the symbolic engine handles the low-level inferential details and guarantees absolute logical correctness [p5, 2, 11].

The neuro-symbolic convergence naturally requires an agentic workflow: Rather than outputting a static answer in a single pass, LLMs are now being deployed as dynamic agents that engage in multi-turn, interactive reasoning loops [p5, 1, 12]. These agents recursively decompose complex mathematical statements into dependency graphs, retrieve necessary axioms from formal libraries, and actively write code that is tested against a compiler [p5, 13, 14]. If a proposed proof step contains a syntax error or a logical fallacy, the formal environment (like the Lean or Coq proof assistant) immediately rejects it and returns a highly specific error message. The AI agent then ingests this diagnostic feedback, reflects on its mistake, and iteratively synthesizes a corrected proof step [p5, 15-17].

Major industry players are investing heavily in this transition, pushing models away from merely "guessing" text toward actively verifying their logic to discover new mathematics:

\begin{itemize}
    \item Google DeepMind recently achieved a historic milestone with its AlphaProof and AlphaGeometry systems, which successfully solved four out of six problems at the 2024 International Mathematical Olympiad (IMO) [p5, 18]. AlphaGeometry perfectly exemplifies the neuro-symbolic approach: it uses a neural network to navigate the difficult, creative steps of geometric proofs (like drawing auxiliary constructions), while relying on a symbolic engine to systematically deduce the resulting facts [p5, 19]. DeepMind has also leveraged machine learning to guide human intuition in advanced research, helping mathematicians discover new patterns and formulate theorems in knot theory and representation theory [p5, 20].
    \item OpenAI has actively pushed its models toward complex, inference-time reasoning. The o3 model series demonstrated a massive leap in capability on the FrontierMath benchmark—a highly guarded dataset composed entirely of unpublished, research-level mathematics designed to prevent data contamination [p5, 21, 22]. While previous general-purpose models scored below 2\% on this benchmark, the o3 model reportedly achieved a substantially higher score of around 25\%, highlighting the power of autonomous, multi-step logical deduction [p5, 22].
    \item DeepSeek and others are advancing systems like DeepSeek-Prover, which tightly integrates model training with formal verification. Through a process called Reinforcement Learning from Proof Assistant Feedback (RLPAF), the model's generated tactics are executed in real-time within the Lean environment. The compiler provides dense, step-by-step validation that serves as a reward signal, explicitly training the neural network to align its outputs with strict logical validity rather than textual mimicry [p5, 23].
\end{itemize}

Ultimately, this industry push aims to transform AI from a tool that merely formalizes existing literature into an autonomous engine for mathematical discovery. In the future, agentic AI systems hooked to formal environments could allow mathematicians to study hundreds of complex equations simultaneously—working out an argument for one equation and deploying the AI to autonomously adapt, formalize, and rigorously verify the proofs for large families of related problems at a scale previously unimaginable [p5, 24].

Despite the rapid advancements by industry leaders in training foundational models for mathematics, achieving reliable, autonomous autoformalization remains a formidable challenge. Current architectures are severely constrained by several fundamental bottlenecks. Primary among these is data scarcity; unlike natural language or standard programming languages, the corpus of aligned, high-quality parallel data mapping informal LaTeX mathematics to formalized Lean4 code is exceedingly small. Furthermore, the translation process is intrinsically non-trivial. Human-written mathematical proofs are inherently underspecified, relying heavily on the reader's contextual intuition and domain expertise to bridge implicit logical gaps. In contrast, formal theorem provers like Lean4 demand exhaustive, explicit structural rigor. Consequently, when large language models attempt direct, single-pass translations, they frequently suffer from cascading logical errors, hallucination, or compilation failure. To mitigate these limitations, recent literature suggests transitioning away from monolithic generation toward structured agentic workflows. By integrating iterative self-training methodologies—where compiler feedback is utilized to continuously refine the model—and establishing human-in-the-loop (HITL) protocols to resolve edge cases, the formalization of complex, graduate-level mathematics becomes highly tractable. Testing these orchestration techniques on specific, localized academic curricula presents a vital opportunity to evaluate the practical viability of these systems.

\section{The Research Question}

The intersection of immense technological potential and persistent translational bottlenecks directly motivates the central inquiry of this project. Specifically, this work seeks to address the following research question: \textbf{To what extent can neuro-symbolic and agentic AI architectures accurately translate informal mathematical proofs (LaTeX) into formal Lean4 code within the domains of algebra and number theory?}
